\begin{figure}[htbp]
\centering
\begin{tikzpicture}
\begin{axis}[
  width=0.8\textwidth, height=6cm,
  xlabel={Radius [mm]}, ylabel={Temperature [K]},
  xmin=0, xmax=5.5, ymin=550, ymax=1450,
  grid=major, grid style={cGrid!60,very thin},
  label style={font=\small}, tick label style={font=\scriptsize},
]
\addplot[color=cFuel,thick,mark=*,mark size=1.1pt] coordinates {
 (0.000,1364.72) (0.177,1363.20) (0.354,1359.85) (0.531,1354.77)
 (0.709,1347.98) (0.886,1339.50) (1.063,1317.51) (1.240,1304.03)
 (1.417,1288.93) (1.594,1272.23) (1.772,1253.94) (1.949,1234.10)
 (2.126,1189.82) (2.303,1165.44) (2.480,1139.58) (2.657,1112.27)
 (2.835,1083.52) (3.012,1021.70) (3.189,988.64) (3.366,954.10)
 (3.543,918.03) (3.720,880.42) (3.898,841.25) (4.075,758.22)
 (4.252,714.04) (4.429,667.23) (4.606,633.91) (4.783,626.42)
 (4.961,615.55) (5.138,608.53) (5.315,601.68)};
\draw[black!45,dashed] (axis cs:4.5,550) -- (axis cs:4.5,1450);
\node[font=\tiny,anchor=south west,text=black!65,rotate=90]
  at (axis cs:4.52,700) {gap / cladding};
\node[font=\tiny,anchor=north west,text=black!65] at (axis cs:0.15,1340)
  {pellet};
\end{axis}
\end{tikzpicture}
\caption[Radial temperature profile]{Radial temperature profile at the midplane
at end of cycle. The drop of about \SI{760}{\kelvin} between the pellet centre
(\SI{1365}{\kelvin}) and the cladding outer surface (\SI{602}{\kelvin}) is
dominated by the low conductivity of the fuel and by the thermal resistance of
the gap.}
\label{fig:radial}
\end{figure}

\subsection{Summary of safety margins}

Table~\ref{tab:marges} gives the final state as evaluated by block~D1.

\begin{table}[htbp]
\centering
\caption[Safety margins at end of cycle]{Safety margins at end of cycle, after
the corrections of section~\ref{sec:seuils}. The thresholds remain indicative
and their validation with the supervisor is still pending.}
\label{tab:marges}
\begin{tabularx}{0.95\textwidth}{@{}Xrrc@{}}
\toprule
\textbf{Criterion} & \textbf{Value} & \textbf{Limit} & \textbf{Status}\\
\midrule
Pellet centreline temperature & \SI{1467}{\kelvin} & \SI{3113}{\kelvin}
  & safe (\SI{47}{\percent})\\
Cladding plastic strain (PCMI) & \num{0} &
  \SI{1.00}{\percent} & safe (\SI{0}{\percent})\\
Cladding creep strain & \SI{0.36}{\percent} &
  \SI{1.00}{\percent} & safe (\SI{36}{\percent})\\
Cladding hoop stress & \SI{79.1}{\mega\pascal} &
  \SI{250}{\mega\pascal} & safe (\SI{32}{\percent})\\
Pellet-cladding gap & \SI{-8.0}{\micro\metre} & \SI{5.0}{\micro\metre}
  & \textcolor{cAlert}{\textbf{exceeded}}\\
\bottomrule
\end{tabularx}
\end{table}

\noindent
The stress limit here is not a constant: it is read from the \texttt{sigmaY}
field computed by the solver, and the value reported is the one at the most
highly loaded point (section~\ref{sec:seuils}).

The rod therefore stays comfortably within the margins on all thermomechanical
criteria, but pellet-cladding contact is established. This result is expected
for a rod at the end of a second cycle and is not alarming in itself: it means
that the operating regime has changed in nature, which is precisely the
information a digital twin should report.

\section{Validation of the surrogate}
\label{sec:validation-emulateur}

\subsection{A misleading first training domain}

The first version of the surrogate had been trained on sixteen simulations
crossing four linear heat rates and four durations between
\SIrange{2000}{6500}{\second}. Its \textit{leave-one-out} validation gave
coefficients of determination above $0.998$ on the three quantities tracked. The
result looked excellent.

It was not. Two observations, made by confronting this domain with the reference
simulation of the previous section, invalidated it.

First, the rod reaches \textbf{thermal equilibrium within a few thousand
seconds}. At the two longest durations of the plan, the temperatures differ by
only a few tenths of a kelvin (table~\ref{tab:saturation}): these levels are
duplicates. Of the four duration values, only two carried information, and the
regression problem was in fact almost one-dimensional — power alone sufficed to
predict the output. The high $R^2$ was therefore mostly measuring how easy the
problem was.

\begin{table}[htbp]
\centering
\caption[Thermal saturation of the first domain]{Thermal saturation in the first
training domain: beyond \SI{5000}{\second}, the simulated duration no longer
carries information.}
\label{tab:saturation}
\begin{tabularx}{0.8\textwidth}{@{}Xrrr@{}}
\toprule
\textbf{Linear heat rate} & \textbf{$T$ at \SI{5000}{\second}} &
\textbf{$T$ at \SI{6500}{\second}} & \textbf{Difference}\\
\midrule
\SI{12}{\kilo\watt\per\metre} & \SI{1036.19}{\kelvin} & \SI{1036.24}{\kelvin}
  & \SI{0.05}{\kelvin}\\
\SI{20}{\kilo\watt\per\metre} & \SI{1389.24}{\kelvin} & \SI{1389.35}{\kelvin}
  & \SI{0.11}{\kelvin}\\
\SI{28}{\kilo\watt\per\metre} & \SI{1787.59}{\kelvin} & \SI{1787.76}{\kelvin}
  & \SI{0.17}{\kelvin}\\
\SI{36}{\kilo\watt\per\metre} & \SI{2186.04}{\kelvin} & \SI{2186.20}{\kelvin}
  & \SI{0.16}{\kelvin}\\
\bottomrule
\end{tabularx}
\end{table}

Second and above all, the reference simulation runs to \SI{6.3e7}{\second}, that
is about \textbf{ten thousand times} beyond the upper bound of the training
domain. Yet the phenomena that matter for safety — swelling, creep, gap closure,
PCMI — unfold precisely over these long time scales, entirely absent from the
design of experiments. The surrogate was therefore modelling only the
\emph{initial thermal transient}, while the dashboard's interactive simulator
implied that it was predicting safety margins.

\subsection{Rebuilding over a physically relevant domain}

The design of experiments was redone entirely over realistic irradiation
durations, from \SIrange{7.9e6}{6.3e7}{\second} (about three months to two
years), crossed with five linear heat rates from
\SIrange{12}{36}{\kilo\watt\per\metre}, that is \num{25} complete OFFBEAT
simulations, all of which ran to completion.

This rebuild was made possible only by compiling the solver locally: a two-year
simulation costs \SI{36}{\second} of computation, so that the complete plan fits
in about thirty minutes.

The gain is visible in the data themselves. Over the new domain, the minimum gap
width \textbf{changes sign}, going from \SI{+32}{\micro\metre} to
\SI{-13}{\micro\metre} depending on the conditions: the surrogate can now learn
gap closure, and therefore the onset of PCMI. The circumferential strain also
becomes non-monotonic and changes sign, reflecting the inward creep of the
cladding before the pellet pushes it back.

Figure~\ref{fig:parite} compares the values predicted by \textit{leave-one-out}
cross-validation with those computed by OFFBEAT over this new domain.

\begin{figure}[htbp]
\centering
\begin{tikzpicture}
\begin{axis}[
  width=0.62\textwidth, height=6.4cm,
  xlabel={Maximum temperature computed by OFFBEAT [K]},
  ylabel={Surrogate prediction [K]},
  xmin=600, xmax=2300, ymin=600, ymax=2300,
  grid=major, grid style={cGrid!60,very thin},
  label style={font=\small}, tick label style={font=\scriptsize},
  legend pos=north west, legend style={font=\scriptsize,draw=black!30},
]
\addplot[only marks,mark=*,mark size=1.9pt,color=cFuel] coordinates {
 (970.82,988.19) (940.09,943.55) (930.02,937.54) (949.74,934.82)
 (936.31,959.21) (1186.41,1189.97) (1139.58,1132.33) (1120.10,1097.81)
 (1142.00,1158.28) (1160.16,1146.84) (1401.84,1387.86) (1328.92,1322.66)
 (1295.66,1313.92) (1405.90,1402.35) (1417.85,1424.59) (1614.18,1597.11)
 (1515.43,1563.43) (1556.00,1542.81) (1717.76,1710.48) (1733.75,1741.45)
 (1845.99,1895.67) (1830.40,1767.51) (1848.80,1895.07) (2061.84,2025.93)
 (2079.84,2075.28)};
\addlegendentry{25 simulations}
\addplot[black!55,dashed,thick,domain=850:2200] {x};
\addlegendentry{$y=x$}
\end{axis}
\end{tikzpicture}
\caption[Surrogate parity plot]{Parity plot of the surrogate for the maximum
temperature, in \textit{leave-one-out} validation over the realistic irradiation
domain. Each point is predicted by a model refitted on the other twenty-four.
The coefficient of determination is $R^2 = 0.9951$.}
\label{fig:parite}
\end{figure}

Table~\ref{tab:surrogate} compares the quality obtained over the two domains.

\begin{table}[htbp]
\centering
\caption[Surrogate quality over both domains]{Surrogate quality in
\textit{leave-one-out} validation over the two training domains. The apparent
degradation reflects a problem that has become genuinely two-dimensional and
non-monotonic, not a regression of the method.}
\label{tab:surrogate}
\begin{tabularx}{\textwidth}{@{}Xccc@{}}
\toprule
\textbf{Predicted quantity} &
\textbf{\begin{tabular}{@{}c@{}}$R^2$ (LOO)\\ short domain\\
 (\SIrange{2}{6.5}{\kilo\second})\end{tabular}} &
\textbf{\begin{tabular}{@{}c@{}}$R^2$ (LOO)\\ realistic domain\\
 (3 months to 2 years)\end{tabular}} &
\textbf{\begin{tabular}{@{}c@{}}Mean\\ error\end{tabular}}\\
\midrule
Maximum temperature & $0.9994$ & $0.9951$ & \SI{19}{\kelvin}\\
Maximum creep strain & --- & $0.9905$ & \num{1.5e-4}\\
Minimum gap width & $0.9992$ & $0.9715$ & \SI{1.8}{\micro\metre}\\
\bottomrule
\end{tabularx}
\end{table}

The coefficients fall, and that is the expected result. Over the short domain
the surrogate was fitting a monotonic relation with a single effective variable;
over the realistic domain it must reproduce non-monotonic quantities that change
sign. An $R^2$ of $0.9951$ on a hard problem is worth more than $0.9994$ on a
trivial one. This episode illustrates a general methodological point: \textbf{a
validation metric says nothing about the relevance of the domain over which it
is computed}.

\subsection{Validation on the safety question}

The decisive test does not bear on a coefficient of determination, but on the
surrogate's ability to answer the question the digital twin is meant to address:
\emph{when does pellet-cladding contact occur?}

The surrogate was therefore queried at \SI{25}{\kilo\watt\per\metre}, a power
\textbf{absent from its training set} (which contains only 12, 18, 24, 30 and
\SI{36}{\kilo\watt\per\metre}), and its prediction compared with the reference
simulation of the previous section (figure~\ref{fig:pcmi}).
